%% ma: L10-B13-0098
%% lop: 10
%% chuong: 5
%% bai: 13
%% dang: 10.13.14
%% loai: DS
%% muc: [TH, TH, TH, VD]
%% dap_an: "SĐĐS"
%% nguon: "chuyen-de-cac-so-dac-trung-cua-mau-so-lieu-khong-ghep-nhom-toan-10-chuong-trinh-moi.pdf (GV. Phan Nhật Linh) (P03), Bài 02 Dạng 2 Phần II Câu 2"
%% kien_thuc: "Lớp 10 Bài 13 (số trung bình, trung vị, tứ phân vị, mốt)"
%% trang_thai: chinh-thuc
%% ghi_chu: "Tự duyệt theo yêu cầu giáo viên 22/10/2026."
\begin{ex}
Hai phương pháp học Toán khác nhau được áp dụng cho lớp 10A và 10B. Sau một học kì, điểm khảo sát của hai lớp như sau:
\begin{center}
\begin{tabular}{|c|c|c|c|c|}\hline
\multicolumn{5}{|c|}{Lớp 10A}\\\hline
$2$ & $7$ & $6$ & $3$ & $9$\\\hline
$8$ & $6$ & $7$ & $9$ & $2$\\\hline
$5$ & $7$ & $5$ & $9$ & $8$\\\hline
$8$ & $7$ & $4$ & $3$ & $5$\\\hline
$5$ & $4$ & $5$ & $7$ & $7$\\\hline
\end{tabular}\qquad
\begin{tabular}{|c|c|c|c|c|}\hline
\multicolumn{5}{|c|}{Lớp 10B}\\\hline
$6$ & $7$ & $6$ & $4$ & $7$\\\hline
$9$ & $3$ & $8$ & $7$ & $5$\\\hline
$5$ & $6$ & $8$ & $7$ & $4$\\\hline
$5$ & $3$ & $10$ & $7$ & $9$\\\hline
$6$ & $7$ & $6$ & $7$ & $5$\\\hline
\end{tabular}
\end{center}
Xét tính đúng sai của các khẳng định sau:
\choiceTF{Điểm trung bình của hai lớp bằng nhau}{\True Trung vị của hai mẫu số liệu bằng nhau}{\True Mốt của hai mẫu số liệu bằng nhau}{Căn cứ vào các số đặc trưng, phương pháp áp dụng ở lớp 10A hiệu quả hơn lớp 10B}
\loigiai{a) Sai: $\overline x_A=\dfrac{148}{25}=5{,}92$; $\overline x_B=\dfrac{157}{25}=6{,}28$.

b) Đúng: mỗi lớp có $25$ giá trị, giá trị thứ $13$ sau khi sắp xếp đều là $6$.

c) Đúng: điểm $7$ xuất hiện nhiều nhất ở cả hai lớp ($6$ lần ở 10A, $7$ lần ở 10B), nên mốt đều là $7$.

d) Sai: lớp 10B có điểm trung bình cao hơn; các tứ phân vị lớp 10A là $4{,}5;6;7{,}5$, lớp 10B là $5;6;7$ (điểm 10B đồng đều hơn). Không thể kết luận phương pháp ở lớp 10A hiệu quả hơn.}
\end{ex}
